This proof is based on the formal definition of \emf, a recent formalization
of an \fstar{} subset~\cite{dm4free}.
%x
We use the same notation and rule names from there, but the proof is
nevertheless quite direct with just minimum familiarity with \emf.
%
We use $E$ to represent \fstar contexts:

\[
\begin{array}{rclcl}
    E &=&    \cdot  & & \\
      & \mid & t E &                 \mid & E t \\
      & \mid & \lambda(x:t). E &     \mid & \lambda(x:E). t \\
      & \mid & (x:t) \rightarrow E & \mid & (x:E) \rightarrow t \\
      & \mid & x:t\{E\} &            \mid & x:E\{t\} \\
      & ~ & \ldots
\end{array}
\]

Contexts present a set of bound variables to their hole. We denote
those variables by $\gamma(E)$:

\[
\begin{array}{lcl}
    \gamma(\cdot) &=& \emptyset \\
    \gamma(\lambda(x:t). E) &=& \{x\} \cup \gamma(E) \\
    \gamma(\lambda(x:E). t) &=& \gamma(E) \\
    \ldots
\end{array}
\]
    % γ(∘) = {}
    % γ(λ(x:t). E) = (x:t) ++ γ(E)
    % γ(λ(x:E). t) = γ(E)
    % ...

We say a context $E$ has type $t_1 \Rightarrow t_2$ for $\Gamma$ (noted
by $\Gamma \vdash E : t_1 \Rightarrow t_2$) when for all $e$ such that
$\Gamma, \gamma(E) \vdash e : t_1$ we also have $\Gamma \vdash E[e] : t_2$.

\subsection{Soundness of splitting the proof obligation}

The proof follows mainly from the following lemma:

\begin{lemma}
    \[
    \infer{\Gamma \vDash E[t_1] = E[t_2]}
          {\Gamma, \gamma(E) \vDash t_1 = t_2}
    \]
\end{lemma}
\begin{proof}
    The proof follows from applying functional extensionality and using
    a carefully crafted abstraction.

    Let $S = \lambda f. E[f (\gamma(E))]$ (where $f (\gamma(E))$
    represents $f$ applied to every variable in $\gamma(E)$,
    sequentially). Also, $\lambda \gamma(E). t$ represents
    abstracting $E$ over the variables in $\gamma(E)$, and
    similarly for $\forall \gamma(E). t$.

    Then,
    \[
        \infer[\mbox{Reduction}]{\Gamma \vDash E[t_1] = E[t_2]}
        {\infer[\mbox{V-EqP}]{\Gamma \vDash S (\lambda \gamma(E). t_1) = S (\lambda \gamma(E). t_2)}
        {\infer[\mbox{V-Ext}]{\Gamma \vDash (\lambda \gamma(E). t_1) = (\lambda \gamma(E). t_2)}
        {\infer[\mbox{V-$\forall$i}]{\Gamma \vDash \forall \gamma(E). t_1 = t_2}
                             {\Gamma, \gamma(E) \vDash t_1 = t_2}}}}
    \]

    Note that the particular set of contexts E does not influence
    the proof.



        % Γ, γ(E) |= t1 = t2
        % ------------------------- [V-∀i, repeatedly]
        % Γ |= forall γ(E). t1 = t2
        % ------------------------------ [V-Ext, repeatedly]
        % Γ |= (λγ(E). t1) = (λγ(E). t2)
        % --------------------------------- [V-EqP, with P = S]
        % Γ |= S (λy(E). t1) = S (γ(E). t2)
        % --------------------------------- [Red]
        % Γ |= E[t1] = E[t2]
\end{proof}

Having such lemma, the theorem is proven as follows:
\[
    \infer{\Gamma \vDash E[\phi]}
         {\Gamma \vDash E[\top]
         & \infer[\mbox{Lemma}]{\Gamma \vDash E[\phi] = E[\top]}{
         \infer[\mbox{PropExt}]{\Gamma, \gamma(E) \vDash \phi = \top}{
         \infer[\mbox{Trivial}]{\Gamma, \gamma(E) \vDash \phi \iff \top}
                               {\Gamma, \gamma(E) \vDash \phi}}}
         }
\]


\subsection{Partial completeness of splitting the proof obligation}

While the previous theorem allows to soundly split any
subformulae within a VC, we have seen that some of them
constraint the system.
%
Here we prove that for a particular set of contexts, splitting
does not make the judgments stronger, and via Theorem 1, then they are
equivalent.

We define a particular shape of ``positive'' contexts $P$.
We don't attempt to follow all of EMF*'s syntax, just the parts
we commonly see as VCs in practice.

\[
\begin{array}{rclcl}
    P &=&    \cdot  & & \\
       & \mid & \phi \land P            &\mid& P \land \phi \\
       & \mid & \forall (x:t). P \\
\end{array}
\]

In \emf, an implication $a \Rightarrow b$ is just sugar for
$\forall (\_:a). b$, so they are considered as well.

% This is sufficient to catch most VCs we encounter. Most likely, a reason
% is that disjunction/existentials can never appear in strictly-positive
% positions due to the inherent conjunctivity of WPs, but we haven't
% proved it. In any case, this property is not crucial in any way.

\begin{theorem}
    Take $\Gamma$, $E$, and $\phi$ such that:
    (1) $\phi$ is squashed,
    (2) $\Gamma \vdash E : prop \rightarrow prop$
    and (3) $\Gamma \vdash \phi : prop$.
    Then the following holds

    \[
        \infer{\Gamma, \gamma(P) \vDash \phi \qquad \Gamma \vDash P[\top]}
              {\Gamma \vDash P[\phi]}
    \]
\end{theorem}
\begin{proof}
    By induction on the structure of $P$, keeping $\Gamma$ and $\phi$
    universally quantified.
    \begin{itemize}
        \item For $P = \cdot$, trivial
        \item For $P = \psi \land P'$, our hypothesis is
              $\Gamma \vDash \psi \land P'[\phi]$.
              %
              Hence,
              $\Gamma \vDash \psi$ and
              $\Gamma \vDash P'[\phi]$.
              %
              By the IH, we get
              $\Gamma, \gamma(P') \vDash \phi$
              (note that $\gamma(P) = \gamma(P')$)
              and
              $\Gamma \vDash P'[\top]$.
              %
              By weakening and conjunction, 
              we can prove $\Gamma, \gamma(P') \vDash \psi \land \phi$
              and $\Gamma \vDash \psi \land P'[\top]$, and conclude.

        \item For $P = P' \land \psi$, the reasoning is analogous.

        \item For $P = \forall (x:t). P'$, our hypothesis is
              $\Gamma \vDash \forall (x:t). P'[\phi]$. Hence,
              we get $\Gamma, x:t \vDash P'[\phi]$ by
              eliminating the quantifier.
              %
              From the IH, we get
              $\Gamma, x:t, \gamma(P') \vDash \phi$
              and
              $\Gamma, x:t \vDash E[\top]$.
              %
              From this last one, we can use V-$\forall$i to
              obtain $\Gamma \vDash \forall (x:t). E[\top]$, and conclude
              (noting that $\gamma(P) = x:t, \gamma(P')$).
    \end{itemize}
\end{proof}
